\section{Posicionamiento de la entrevista: no es publicidad corporativa, sino una narrativa de paradigma de cómputo}

El valor de esta entrevista de 63 minutos no radica en los resultados trimestrales ni en las especificaciones de un chip en particular, sino en el relato continuo de Jensen Huang sobre la «migración del paradigma de cómputo»: del cómputo paralelo gráfico al cómputo paralelo de propósito general, y de ahí a la reestructuración del software y la industria impulsada por la IA.

La estrategia de preguntas de la conductora Cleo Abram también es clara: omite los temas financieros y del mercado de capitales, y va directo a tres niveles.
\begin{itemize}
    \item Pasado: por qué NVIDIA pasó de ser una empresa de gráficos para videojuegos a convertirse en la empresa central de infraestructura de IA.
    \item Presente: por qué el aprendizaje profundo se convirtió en el paradigma dominante y cómo transformó la pila de software.
    \item Futuro: la Physical AI, la robótica, las restricciones de consumo energético y las decisiones personales.
\end{itemize}

\begin{importantbox}{Planteamiento inicial}
Jensen ofrece, desde el inicio, un juicio con alta confianza: \textit{``Everything that moves will be robotic someday, and it will be soon.''}

El sentido de esta frase no es que «todos los dispositivos se vuelvan humanoid de inmediato», sino que la capacidad de automatización «que puede percibir, decidir y ejecutar» penetrará en los sistemas móviles e industriales, con una tendencia de aceleración en los próximos años.
\end{importantbox}

\begin{knowledgebox}{Densidad técnica de la entrevista}
Aunque el video está dirigido al público general, en realidad abarca varios niveles de temas técnicos:
\begin{itemize}
    \item Arquitectura de cómputo: la diferenciación de roles entre el CPU serial y el GPU parallel.
    \item Modelo de programación: cómo CUDA reduce la barrera de entrada al cómputo paralelo.
    \item Paradigma de aprendizaje: el Software 2.0 después de AlexNet.
    \item Entrenamiento y despliegue: del centro de datos a la AI supercomputer personal.
    \item Infraestructura robótica: la lógica de combinación de Omniverse + Cosmos.
    \item Seguridad y gobernanza: descomponer la AI safety en problemas ejecutables de ingeniería.
\end{itemize}
\end{knowledgebox}

\subsection{Resumen de la sección}
Esta entrevista puede considerarse como la explicación pública de la «estrategia de plataforma» de NVIDIA: el núcleo no es apostar por alguna aplicación de moda, sino construir una plataforma de cómputo de propósito general capaz de sostener, de manera duradera, la siguiente generación de algoritmos y su implementación industrial.

